% Example "oscillator": the problem, the solutions against the exact one, and
% the energy, phase and stability tests. Every number comes from data/
% (parameters.tex, results.tex, exact_tail.dat, tail_<method>.dat,
% portrait_<method>.dat, energy_<method>.dat, theory_<method>.dat), written by
% examples/oscillator/main.f90: run `fpm run --example oscillator` from the
% project root, then build this file in examples/oscillator/.
\documentclass[11pt,a4paper]{article}
\usepackage[margin=2.2cm]{geometry}
\usepackage{amsmath,booktabs}
\usepackage{pgfplots}
\usepgfplotslibrary{groupplots}
\pgfplotsset{compat=1.18}
\pagestyle{empty}

\begin{document}
\section*{Example: harmonic oscillator}

\subsection*{1. Problem}
\begin{equation*}
  \frac{d^2x}{dt^2} = -\omega^2 x
  \quad\Longleftrightarrow\quad
  \frac{d}{dt}\begin{pmatrix} x \\ v \end{pmatrix}
  = \begin{pmatrix} 0 & 1 \\ -\omega^2 & 0 \end{pmatrix}
    \begin{pmatrix} x \\ v \end{pmatrix},
  \qquad x(0) = 1,\; v(0) = 0,
\end{equation*}
\begin{equation*}
  \text{exact solution}\qquad x(t) = \cos\omega t, \quad v(t) = -\omega\sin\omega t,
  \qquad E = \tfrac12\bigl(v^2 + \omega^2 x^2\bigr) = E_0 \text{ for all } t .
\end{equation*}
The matrix has eigenvalues $\pm i\omega$: the exact flow is a rotation that
keeps the energy forever. One Runge--Kutta step multiplies the solution by
the stability polynomial $R(i\omega h)$, a rotation by $\arg R$ combined with
a scaling by $|R|$ (computed from the Butcher tableau alone).
\begin{center}
\small
% Written by examples/oscillator/main.f90: do not edit.
\begin{tabular}{ll}
\toprule
frequency & $\omega = 1$ (period $2\pi$) \\
initial condition & $x(0) = 1$, $v(0) = 0$, $E_0 = 1/2$ \\
step & $h = \Delta t = 0.1$ \\
interval & 100 periods, 6283 steps, $t_{\mathrm{end}} = 628.3$ \\
methods & \texttt{rk3ws}, \texttt{rk4}, \texttt{dp54} \\
interface & plain arrays, \texttt{rk\_array\_scheme\_t} \\
figures & $x(t)$ over the last 3 periods; $(x, v)$ every 4 steps \\
phase measured & $\operatorname{atan2}(-v, x) - t$, wrapped to $(-\pi, \pi]$ \\
limit from $R$ & first $y$ with $|R(iy)| > 1$: scan by $10^{-3}$, then 60 bisections \\
limit measured & largest $h$ with $E_N \le E_0$ after $N = 2000$ steps: 50 bisections on $[0, 6]$ \\
\bottomrule
\end{tabular}

\end{center}

\subsection*{2. Solutions}
\begin{center}
\begin{tikzpicture}
  \begin{groupplot}[
      group style={group size=2 by 1, horizontal sep=1.8cm}, % two panels side by side
      width=7.6cm, height=6.2cm,                           % size of each panel
      grid=major, grid style={black!12},                   % light grid
    ]
    % ---- (a) x(t) over the last periods ----
    \nextgroupplot[
      title={(a) $x(t)$, last 3 of 100 periods},           % panel title
      xlabel={$t$}, ylabel={$x$}, ymin=-1.2, ymax=1.2,      % position against time
      legend to name=osclegend, legend columns=5,          % legend below both panels
      legend style={font=\small, /tikz/every even column/.append style={column sep=0.3cm}},
    ]
    \addplot[black, thick] table[x=t, y=x] {data/exact_tail.dat};                % exact cos t
    \addlegendentry{exact}
    \addplot[green!50!black, thick] table[x=t, y=x] {data/tail_rk3ws.dat};      % order 3: smaller
    \addlegendentry{order 3}
    \addplot[blue!80!black, thick, dashed] table[x=t, y=x] {data/tail_rk4.dat}; % order 4: on the exact curve
    \addlegendentry{\texttt{rk4}}
    \addplot[red!75!black, thick, dotted] table[x=t, y=x] {data/tail_dp54.dat}; % order 5: on the exact curve
    \addlegendentry{\texttt{dp54}}
    % ---- (b) phase portrait over 100 periods ----
    \nextgroupplot[
      title={(b) $(x, v)$ over 100 periods},               % panel title
      xlabel={$x$}, ylabel={$v$},                          % phase plane
      axis equal image, xmin=-1.15, xmax=1.15, ymin=-1.15, ymax=1.15, % circles stay circles
    ]
    \addplot[green!50!black, only marks, mark=*, mark size=0.3pt] table[x=x, y=v] {data/portrait_rk3ws.dat};  % order 3: spirals in
    \addplot[blue!80!black, only marks, mark=*, mark size=0.3pt] table[x=x, y=v] {data/portrait_rk4.dat};    % order 4: stays on the circle
    \addplot[black, thick, domain=0:360, samples=181] ({cos(x)}, {sin(x)}); % exact orbit: the unit circle
  \end{groupplot}
  \node[below=0.2cm] at (current bounding box.south)        % legend centred under both panels
    {\pgfplotslegendfromname{osclegend}};
\end{tikzpicture}
\end{center}
(a) After 100 periods \texttt{rk3ws}, of order 3 on this linear problem, is
$2.6\,\%$ too small in amplitude and nearly in phase; \texttt{rk4} and
\texttt{dp54} are on the exact curve at this scale. (b) The same in the phase
plane: \texttt{rk3ws} spirals inwards, \texttt{rk4} stays on the circle.

\newpage
\subsection*{3. Test: energy, phase and the imaginary-axis limit}

\paragraph{Prediction.} After $n$ steps, $E/E_0 = |R(ih)|^{2n}$ and the
phase error is $n\,(\arg R(ih) - h)$. Methods of order 1 and 2 have
$|R(iy)| > 1$ for every $y > 0$: the energy grows at any step. Orders 3 and 4
are stable on the imaginary axis up to $\sqrt 3$ and $2\sqrt 2$.

\paragraph{Result.}
\begin{center}
\small
% Written by examples/oscillator/main.f90: do not edit.
\begin{tabular}{lcccccc}
\toprule
& \multicolumn{2}{c}{$E/E_0$ after 100 periods} & \multicolumn{2}{c}{phase error (rad)} & \multicolumn{2}{c}{imaginary-axis limit} \\
\cmidrule(lr){2-3}\cmidrule(lr){4-5}\cmidrule(lr){6-7}
name & predicted & measured & predicted & measured & from $R$ & measured \\
\midrule
\texttt{rk3ws} & $0.949154$ & $0.949154$ & $2.09\times10^{-3}$ & $2.09\times10^{-3}$ & 1.7321 & 1.7321 \\
\texttt{rk4} & $0.999913$ & $0.999913$ & $-5.22\times10^{-4}$ & $-5.22\times10^{-4}$ & 2.8284 & 2.8284 \\
\texttt{dp54} & $0.999997$ & $0.999997$ & $2.98\times10^{-7}$ & $2.98\times10^{-7}$ & 0.9972 & 0.9972 \\
\bottomrule
\end{tabular}

\end{center}

\begin{center}
\begin{tikzpicture}
  \begin{loglogaxis}[
      width=10.5cm, height=7cm,                           % figure size
      xlabel={time $t$},                                   % x: time, 1 step to 100 periods
      ylabel={energy error $|E/E_0 - 1|$},                 % y: relative energy drift
      ymin=1e-12, ymax=1,                                  % from RK4's drift to order one
      grid=major, grid style={black!12},                   % light grid
      legend pos=outer north east, legend cell align=left, % legend right of the axes
      legend style={font=\small},                          % compact legend
    ]
    \addplot[green!50!black, thick]  table[x=t, y=drift] {data/theory_rk3ws.dat};    % order 3, predicted
    \addlegendentry{\texttt{rk3ws}}
    \addplot[blue!80!black, thick]   table[x=t, y=drift] {data/theory_rk4.dat};      % RK4, predicted
    \addlegendentry{\texttt{rk4}}
    \addplot[red!75!black, thick]    table[x=t, y=drift] {data/theory_dp54.dat};     % DP5(4), predicted
    \addlegendentry{\texttt{dp54}}
    \addplot[only marks, mark=*, mark size=1.2pt, green!50!black]  table[x=t, y=drift] {data/energy_rk3ws.dat}; % order 3, measured
    \addplot[only marks, mark=*, mark size=1.2pt, blue!80!black]   table[x=t, y=drift] {data/energy_rk4.dat};   % RK4, measured
    \addplot[only marks, mark=*, mark size=1.2pt, red!75!black]    table[x=t, y=drift] {data/energy_dp54.dat};  % DP5(4), measured
  \end{loglogaxis}
\end{tikzpicture}
\end{center}
Lines: $\bigl||R(ih)|^{2n} - 1\bigr|$; dots: the runs. Explicit Runge--Kutta
methods do not conserve energy: orders 1--2 gain it, orders 3--5 lose it, at
the rate $|R(ih)|^2$ sets (Euler $1 + h^2$, order 2 $1 + h^4/4$, order 3
$1 - h^4/12 + h^6/36$, RK4 $1 - h^6/72 + h^8/576$). The drift is linear in $t$
while it is small, then exponential. \texttt{dp54} is stable on the
imaginary axis only up to $h \approx 0.997$, yet at $h = 0.1$ its energy
error is the smallest: a stability limit and an accuracy at one step are
different questions.
\end{document}
